\documentclass[aspectratio=169]{beamer}

\usetheme{metropolis}
\usepackage[utf8]{inputenc}
\usepackage[T1]{fontenc}
\usepackage[brazil]{babel}
\usepackage{listings}
\usepackage{tabularx}
\usepackage{array}
\usepackage{booktabs}
\newcolumntype{R}[1]{>{\raggedright\arraybackslash}p{#1}}
\usepackage{tikz}
\usetikzlibrary{arrows.meta, positioning, shapes.geometric, calc, fit, backgrounds}

\setbeamertemplate{frame numbering}[fraction]
\metroset{block=fill}

\definecolor{termbg}{RGB}{40,44,52}
\definecolor{termfg}{RGB}{220,220,220}
\definecolor{codekw}{RGB}{235,129,27}
\definecolor{codecom}{RGB}{140,150,160}
\definecolor{codestr}{RGB}{152,195,121}

\lstdefinestyle{terminal}{
    backgroundcolor=\color{termbg},
    basicstyle=\ttfamily\footnotesize\color{termfg},
    breaklines=true,
    frame=single,
    rulecolor=\color{termbg},
    framerule=2pt,
    xleftmargin=4pt,
    xrightmargin=4pt,
    columns=fullflexible,
    keepspaces=true,
    showstringspaces=false,
    literate={á}{{\'a}}1 {à}{{\`a}}1 {â}{{\^a}}1 {ã}{{\~a}}1 {é}{{\'e}}1 {ê}{{\^e}}1
             {í}{{\'i}}1 {ó}{{\'o}}1 {ô}{{\^o}}1 {õ}{{\~o}}1 {ú}{{\'u}}1 {ç}{{\c c}}1
}

\lstdefinestyle{codigo}{
    style=terminal,
    language=C,
    keywordstyle=\color{codekw}\bfseries,
    commentstyle=\color{codecom}\itshape,
    stringstyle=\color{codestr},
    morekeywords={pthread_t, sem_t, size_t},
    moredelim=[is][\color{codekw}\bfseries]{@}{@}
}

\lstdefinestyle{java}{
    style=codigo,
    language=Java,
    morekeywords={synchronized}
}

% Elementos visuais das figuras
\tikzset{
    proc/.style={draw=mDarkTeal, fill=mDarkTeal!8, rounded corners=3pt, thick,
        minimum width=2.1cm, minimum height=0.8cm, align=center, font=\small},
    procf/.style={proc, draw=mDarkTeal!70, fill=white, minimum width=1.3cm, font=\footnotesize},
    destaque/.style={proc, draw=mLightBrown, fill=mLightBrown!12},
    seta/.style={-{Stealth[length=2.2mm]}, thick, draw=mDarkTeal},
    setad/.style={-{Stealth[length=2.2mm]}, very thick, draw=mLightBrown},
    rotulo/.style={font=\scriptsize, fill=white, inner sep=1.5pt},
    celula/.style={draw=mDarkTeal, minimum width=1.05cm, minimum height=0.8cm, font=\ttfamily\small, fill=white},
    cheia/.style={celula, fill=mLightBrown!25}
}

% Revelar partes de uma figura TikZ aos poucos (overlays do beamer).
\tikzset{
    invisivel/.style={opacity=0, text opacity=0},
    visible on/.style={alt={#1{}{invisivel}}},
    alt/.code args={<#1>#2#3}{\alt<#1>{\pgfkeysalso{#2}}{\pgfkeysalso{#3}}}
}

\title{Concorrência e sincronização}
\subtitle{Condição de corrida, semáforos, produtor-consumidor e monitores}
\author{Weslei Ferreira Santos}
\institute{MATA58 · Sistemas Operacionais\\Instituto de Computação / UFBA}
\date{7 de outubro de 2026}

\begin{document}

\maketitle

\begin{frame}{Agenda}
\begin{columns}[T]
\column{0.5\textwidth}
\textbf{Parte 1: teoria} \hfill {\small 45 min}
\begin{enumerate}
    \item De processos a threads
    \item Condição de corrida
    \item Região crítica e Peterson
    \item Semáforos
    \item Produtor-consumidor em C
    \item Ver funcionando: simulação
    \item Por que monitores?
\end{enumerate}

\column{0.5\textwidth}
\textbf{Parte 2: prática} \hfill {\small 45 min}
\begin{enumerate}
    \setcounter{enumi}{7}
    \item Seis lacunas no produtor-consumidor (C)
    \item Mailbox sem sincronização (Java, 3b)
    \item Mailbox com monitor (Java, 3s)
    \item Atividade final
\end{enumerate}
\end{columns}
\end{frame}

\begin{frame}{A pergunta desta aula}
\begin{center}
{\Large \textbf{Se as threads compartilham a memória,\\[0.2em] como impedir que uma atrapalhe a outra?}}
\end{center}

\vspace{1em}
\textbf{Ao final, você deve conseguir:}
\begin{itemize}
    \item explicar por que \texttt{contador++} em duas threads perde incrementos;
    \item definir região crítica e exclusão mútua;
    \item dizer o que \texttt{mutex}, \texttt{empty} e \texttt{full} controlam, e por que a ordem importa;
    \item completar o produtor-consumidor com os semáforos certos;
    \item achar a condição de corrida na Mailbox 3b e explicar como \texttt{synchronized}, \texttt{wait} e \texttt{notify} a resolvem na 3s.
\end{itemize}
\end{frame}

% ============================================================
\section{De processos a threads}
% ============================================================

\begin{frame}{Lembrando a aula passada}
\begin{columns}[c]
\column{0.5\textwidth}
\begin{itemize}
    \item No \texttt{fork2.c}, o pai somou \textbf{0}.
    \item Cada processo tem a \textbf{sua} memória: o que um escreve, o outro não vê.
    \item Para conversar, foi preciso um canal do SO: o \textbf{pipe}.
\end{itemize}

\column{0.48\textwidth}
\begin{alertblock}{Hoje o problema se inverte}
As \textbf{threads} de um processo \textbf{dividem} a memória. Comunicar ficou fácil. O difícil agora é \textbf{não se atrapalhar}.
\end{alertblock}
\end{columns}
\end{frame}

\begin{frame}{O que as threads compartilham}
\begin{columns}[c]
\column{0.5\textwidth}
\begin{center}
\begin{tikzpicture}[node distance=0pt,
    reg/.style={draw=mDarkTeal, fill=mDarkTeal!8, minimum width=4.6cm, minimum height=0.55cm, font=\small},
    pil/.style={draw=mLightBrown, fill=mLightBrown!15, minimum width=2.2cm, minimum height=0.9cm, font=\small, align=center}]
    \node[reg] (cod) {código};
    \node[reg, below=of cod] (glo) {variáveis globais};
    \node[reg, below=of glo] (heap) {heap (\texttt{malloc})};
    \node[reg, below=of heap] (arq) {arquivos abertos};
    \node[pil, below=0.25cm of arq.south west, anchor=north west] (p1) {pilha\\thread 1};
    \node[pil, below=0.25cm of arq.south east, anchor=north east] (p2) {pilha\\thread 2};
    \node[draw=black!40, dashed, rounded corners, fit=(cod)(p1)(p2), inner sep=0.2cm,
          label={[font=\scriptsize, text=black!60]above:um processo}] {};
\end{tikzpicture}
\end{center}

\column{0.48\textwidth}
\small
\textbf{Compartilhado} (azul): código, globais, heap, arquivos.\\[0.5em]
\textbf{De cada thread} (laranja): pilha (variáveis locais) e registradores.

\vspace{0.6em}
\begin{exampleblock}{Consequência}
Uma variável global escrita por uma thread é vista pela outra na hora, sem pipe.
\end{exampleblock}
\end{columns}
\end{frame}

\begin{frame}[fragile]{Criando threads em C}
\begin{lstlisting}[style=codigo]
void *somar(void *arg) { ... return NULL; }   /* o que a thread executa */

pthread_t t1, t2;
pthread_create(&t1, NULL, somar, NULL);        /* cria e ja comeca a rodar */
pthread_create(&t2, NULL, somar, NULL);
pthread_join(t1, NULL);                        /* espera a thread terminar */
pthread_join(t2, NULL);
\end{lstlisting}

\begin{center}
\footnotesize
\begin{tabular}{@{}lll@{}}
\toprule
& \textbf{Processos (aula passada)} & \textbf{Threads (hoje)} \\
\midrule
Criar & \texttt{fork()} & \texttt{pthread\_create()} \\
Esperar terminar & \texttt{waitpid()} & \texttt{pthread\_join()} \\
Memória & cópia separada & \textbf{a mesma} \\
Compilar & normal & com \texttt{-pthread} \\
\bottomrule
\end{tabular}
\end{center}
\end{frame}

\begin{frame}[fragile]{Compilando e executando os exemplos}
\begin{lstlisting}[style=terminal, basicstyle=\ttfamily\scriptsize\color{termfg}]
$ gcc -Wall -Wextra examples/dia07/corrida.c -o corrida -pthread
$ ./corrida
$ javac -d build/java-3b examples/dia07/mailbox-3b/*.java
$ timeout 5 java -cp build/java-3b ThreadSync3b
\end{lstlisting}

\begin{center}
\scriptsize
\begin{tabular}{@{}>{\ttfamily}l>{\raggedright\arraybackslash}p{9.2cm}@{}}
\toprule
\textrm{\textbf{Parte}} & \textbf{O que significa} \\
\midrule
-pthread & Prepara o programa para usar threads e semáforos (use sempre hoje). \\
-o corrida & Nome do executável; \texttt{./corrida} executa. \\
javac -d build/java-3b & Compila todos os \texttt{.java} da pasta (\texttt{*.java}) e guarda os \texttt{.class} em \texttt{build/java-3b}. \\
java -cp build/java-3b & Executa a classe indicada, procurando os \texttt{.class} nessa pasta. \\
timeout 5 & Encerra depois de 5 segundos: as mailboxes nunca terminam sozinhas. \\
\bottomrule
\end{tabular}
\end{center}

\scriptsize Na prática 1, acrescente \texttt{-std=c11 -Wpedantic -Werror}: avisos viram erros.
\end{frame}

% ============================================================
\section{Condição de corrida}
% ============================================================

\begin{frame}{\texttt{contador++} não é uma operação só}
\begin{center}
{\large Para o processador, \texttt{contador++} são \textbf{três passos}:}

\vspace{0.8em}
\begin{tikzpicture}[node distance=1.2cm]
    \node[proc] (l) {1. \textbf{ler}\\\scriptsize valor da memória};
    \node[proc, right=of l] (s) {2. \textbf{somar 1}\\\scriptsize no registrador};
    \node[proc, right=of s] (e) {3. \textbf{escrever}\\\scriptsize de volta na memória};
    \draw[seta] (l) -- (s);
    \draw[seta] (s) -- (e);
\end{tikzpicture}
\end{center}

\vspace{0.6em}
\begin{alertblock}{O perigo}
Entre um passo e outro, o escalonador pode trocar de thread. A outra thread pode ler o valor \textbf{antigo} antes da primeira escrever o novo.
\end{alertblock}
\end{frame}

\begin{frame}{Uma intercalação que perde um incremento}
\small Contador começa em 0. Cada thread quer somar 1. Esperado: \textbf{2}.

\begin{center}
\begin{tabular}{@{}clll@{}}
\toprule
\textbf{Passo} & \textbf{Thread A} & \textbf{Thread B} & \textbf{contador} \\
\midrule
\uncover<2->{1} & \uncover<2->{lê 0} & & \uncover<2->{0} \\
\uncover<3->{2} & & \uncover<3->{lê 0} & \uncover<3->{0} \\
\uncover<4->{3} & \uncover<4->{calcula 1} & & \uncover<4->{0} \\
\uncover<5->{4} & & \uncover<5->{calcula 1} & \uncover<5->{0} \\
\uncover<6->{5} & \uncover<6->{escreve 1} & & \uncover<6->{1} \\
\uncover<7->{6} & & \uncover<7->{escreve 1} & \uncover<7->{\textbf{1}} \\
\bottomrule
\end{tabular}
\end{center}

\uncover<8->{
\begin{alertblock}{Condição de corrida}
O resultado depende da \textbf{ordem} em que as threads foram intercaladas. Aqui, a escrita de A foi apagada pela de B.
\end{alertblock}}
\end{frame}

\begin{frame}[fragile]{Na prática: \texttt{corrida.c}}
\begin{lstlisting}[style=codigo, basicstyle=\ttfamily\scriptsize\color{termfg}]
static volatile int contador = 0;           /* global: as duas threads veem a mesma */

static void *somar(void *arg) {
    for (int i = 0; i < 1000000; ++i)
        contador = contador + 1;            /* regiao critica SEM protecao */
    return NULL;
}
\end{lstlisting}

\begin{lstlisting}[style=terminal, basicstyle=\ttfamily\scriptsize\color{termfg}]
$ gcc -Wall -Wextra examples/dia07/corrida.c -o corrida -pthread
$ ./corrida
Contador: 1006933 (esperado: 2000000)
$ ./corrida
Contador: 1349235 (esperado: 2000000)
\end{lstlisting}

\small Quase metade dos incrementos se perde, e o resultado \textbf{muda a cada execução}.
\end{frame}

\begin{frame}{Dois avisos sobre \texttt{corrida.c}}
\begin{itemize}
    \item Em C, duas threads acessando a mesma variável sem sincronização é uma \textbf{data race}: comportamento \textbf{indefinido}. O resultado não está garantido nem a ser ``1 ou 2''.
    \item \texttt{volatile} \textbf{não conserta}: ele só impede o compilador de juntar as somas. Está ali para a corrida aparecer.
    \item Pôr um \texttt{sleep} também não conserta: só muda a probabilidade.
\end{itemize}

\vspace{0.6em}
\begin{exampleblock}{O que conserta}
Garantir que só uma thread por vez execute o trecho que mexe no contador. Isso tem nome: \textbf{exclusão mútua}.
\end{exampleblock}
\end{frame}

% ============================================================
\section{Região crítica e exclusão mútua}
% ============================================================

\begin{frame}{Região crítica e exclusão mútua}
\begin{alertblock}{Definições}
\textbf{Região crítica}: trecho do código que acessa um recurso compartilhado.\\[0.3em]
\textbf{Exclusão mútua}: no máximo \textbf{uma} thread por vez dentro da região crítica.
\end{alertblock}

\textbf{Uma boa solução precisa garantir} (Tanenbaum):
\begin{enumerate}
    \item Duas threads nunca estão ao mesmo tempo na região crítica.
    \item Não depende da velocidade nem do número de CPUs.
    \item Quem está fora da região crítica não bloqueia ninguém.
    \item Ninguém espera para sempre para entrar.
\end{enumerate}
\end{frame}

\begin{frame}[fragile]{Algoritmo de Peterson (duas threads)}
\begin{columns}[T]
\column{0.55\textwidth}
\begin{lstlisting}[style=terminal, basicstyle=\ttfamily\scriptsize\color{termfg}]
/* entrada */
interessado[eu] = verdadeiro
vez = outro
enquanto interessado[outro] e vez == outro:
    esperar

/* região crítica */

/* saída */
interessado[eu] = falso
\end{lstlisting}

\column{0.43\textwidth}
\small
\begin{itemize}
    \item \texttt{interessado[i]}: a thread \texttt{i} quer entrar.
    \item \texttt{vez}: quem cede em caso de empate.
\end{itemize}

\vspace{0.3em}
\begin{tabular}{@{}R{2.4cm}R{3.2cm}@{}}
\toprule
\textbf{Situação} & \textbf{O que acontece} \\
\midrule
O outro não quer entrar & Entro direto \\
Os dois querem & Quem escreveu \texttt{vez} por \textbf{último} espera \\
\bottomrule
\end{tabular}
\end{columns}
\end{frame}

\begin{frame}[fragile]{Peterson em C: \texttt{peterson-code.c}}
\begin{lstlisting}[style=codigo]
static void enter_region(int eu) {
    int outro = 1 - eu;
    atomic_store(&interested[eu], 1);       /* quero entrar */
    atomic_store(&turn, outro);             /* cedo a vez em caso de empate */
    while (atomic_load(&interested[outro]) && atomic_load(&turn) == outro)
        ;                                   /* espera ocupada */
}
static void leave_region(int eu) { atomic_store(&interested[eu], 0); }
\end{lstlisting}

\begin{lstlisting}[style=terminal, basicstyle=\ttfamily\scriptsize\color{termfg}]
$ gcc -Wall -Wextra examples/dia07/peterson-code.c -o peterson -pthread && ./peterson
Contador: 20000 (esperado: 20000)
\end{lstlisting}

\small \textbf{Por que atomics?} Com variáveis comuns, compilador e processador podem reordenar leituras e escritas, e Peterson deixa de funcionar.
\end{frame}

\begin{frame}{Espera ocupada}
\textbf{O que a thread faz enquanto está presa no \texttt{while}?}

\pause
\begin{itemize}
    \item Fica \textbf{girando}: testa, testa, testa de novo.
    \item Continua \textbf{gastando CPU} sem fazer trabalho útil.
    \item Com uma CPU só, pode até atrasar a thread que precisa sair da região crítica.
\end{itemize}

\vspace{0.6em}
\begin{alertblock}{Queremos melhor}
Em vez de girar, a thread deveria \textbf{dormir} e ser \textbf{acordada} quando puder seguir. É isso que o semáforo faz.
\end{alertblock}
\end{frame}

% ============================================================
\section{Semáforos}
% ============================================================

\begin{frame}{Semáforo: um contador que nunca fica negativo}
\begin{columns}[T]
\column{0.5\textwidth}
\textbf{\texttt{sem\_wait(\&s)}} (também: \emph{down}, P)
\begin{itemize}
    \item Se \texttt{s > 0}: diminui 1 e segue.
    \item Se \texttt{s == 0}: a thread \textbf{dorme}.
\end{itemize}

\column{0.5\textwidth}
\textbf{\texttt{sem\_post(\&s)}} (também: \emph{up}, V)
\begin{itemize}
    \item Soma 1.
    \item Se alguém dorme em \texttt{s}, \textbf{acorda} um.
\end{itemize}
\end{columns}

\vspace{0.8em}
\begin{exampleblock}{Analogia: estacionamento com placa de vagas}
O valor do semáforo é o número de vagas na placa. Quem chega e vê 0 \textbf{espera parado}, sem rodar em volta do quarteirão. Quem sai, soma uma vaga e libera o próximo.
\end{exampleblock}

\small A espera e o despertar são feitos pelo SO de forma \textbf{atômica}: não há corrida no próprio contador.
\end{frame}

\begin{frame}{Os três semáforos do produtor-consumidor}
\begin{center}
\begin{tabular}{@{}llR{7.4cm}@{}}
\toprule
\textbf{Semáforo} & \textbf{Começa em} & \textbf{O que controla} \\
\midrule
\texttt{mutex} & 1 & \textbf{Exclusão mútua}: só uma thread mexe no buffer por vez. É a ``chave''. \\
\texttt{empty} & N & \textbf{Disponibilidade}: quantas \textbf{vagas} livres existem. \\
\texttt{full} & 0 & \textbf{Disponibilidade}: quantos \textbf{itens} prontos existem. \\
\bottomrule
\end{tabular}
\end{center}

\vspace{0.6em}
\begin{alertblock}{Dois problemas diferentes}
\texttt{mutex} responde ``\textbf{posso mexer agora?}''.\\
\texttt{empty} e \texttt{full} respondem ``\textbf{tem o que eu preciso?}''.
\end{alertblock}
\end{frame}

\begin{frame}{Simulando no quadro com N = 3}
\begin{center}
\begin{tabular}{@{}lccc@{}}
\toprule
\textbf{Evento} & \textbf{empty} & \textbf{full} & \textbf{count} \\
\midrule
início & 3 & 0 & 0 \\
\uncover<2->{produz A} & \uncover<2->{2} & \uncover<2->{1} & \uncover<2->{1} \\
\uncover<3->{produz B} & \uncover<3->{1} & \uncover<3->{2} & \uncover<3->{2} \\
\uncover<4->{consome A} & \uncover<4->{2} & \uncover<4->{1} & \uncover<4->{1} \\
\uncover<5->{produz C} & \uncover<5->{1} & \uncover<5->{2} & \uncover<5->{2} \\
\uncover<6->{produz D} & \uncover<6->{0} & \uncover<6->{3} & \uncover<6->{3} \\
\bottomrule
\end{tabular}
\end{center}

\uncover<7->{
\textbf{E se o produtor quiser o quinto item?} Ele dorme em \texttt{sem\_wait(\&empty)}, sem segurar o \texttt{mutex}, até o consumidor liberar uma vaga.}

\uncover<8->{\small Repare: \texttt{empty + full = N} sempre que nenhuma operação está em andamento.}
\end{frame}

\begin{frame}{Sua vez: N = 5}
\begin{center}
{\large Buffer com \textbf{N = 5}: \texttt{empty} começa em 5 e \texttt{full} em 0.}

\vspace{1em}
{\large Depois de \textbf{dois itens produzidos} e \textbf{um consumido},\\[0.2em] quanto valem \texttt{empty} e \texttt{full}?}
\end{center}

\pause
\vspace{1em}
\begin{exampleblock}{Resposta}
\texttt{empty = 4} e \texttt{full = 1}. De novo, \texttt{empty + full = N}.
\end{exampleblock}
\end{frame}

\begin{frame}[fragile]{Voltando ao \texttt{corrida.c}: dois consertos}
\begin{columns}[T]
\column{0.48\textwidth}
\textbf{Operação atômica}
\begin{lstlisting}[style=codigo, basicstyle=\ttfamily\scriptsize\color{termfg}]
static atomic_int contador = 0;
...
atomic_fetch_add(&contador, 1);
\end{lstlisting}
\scriptsize A soma inteira vira uma operação indivisível (\texttt{lock addl} no processador).

\column{0.48\textwidth}
\textbf{Mutex}
\begin{lstlisting}[style=codigo, basicstyle=\ttfamily\scriptsize\color{termfg}]
pthread_mutex_lock(&trava);
contador = contador + 1;
pthread_mutex_unlock(&trava);
\end{lstlisting}
\scriptsize Só uma thread por vez; quem encontra a trava ocupada dorme.
\end{columns}

\begin{lstlisting}[style=terminal, basicstyle=\ttfamily\scriptsize\color{termfg}]
$ gcc -Wall -Wextra examples/dia07/corrida_atomic.c -o corrida-atomic -pthread
$ ./corrida-atomic
Contador: 2000000 (esperado: 2000000)
$ gcc -Wall -Wextra examples/dia07/corrida_mutex.c -o corrida-mutex -pthread
$ ./corrida-mutex
Contador: 2000000 (esperado: 2000000)
\end{lstlisting}
\end{frame}

% ============================================================
\section{Produtor-consumidor em C}
% ============================================================

\begin{frame}{O problema: um buffer, dois ritmos}
\begin{center}
\begin{tikzpicture}
    \node[proc] (prod) at (-1.2,0) {Produtor};
    \foreach \i/\v in {0/7, 1/42, 2/13, 3/{}, 4/{}} {
        \ifx\v\empty \node[celula] (c\i) at (2.0+\i*1.05,0) {}; \else \node[cheia] (c\i) at (2.0+\i*1.05,0) {\v}; \fi
        \node[font=\tiny, below=1pt of c\i] {\i};
    }
    \node[proc] (cons) at (9.9,0) {Consumidor};
    \draw[setad] (prod) -- (c0.west);
    \draw[setad] (c4.east) -- (cons);
    \node[font=\scriptsize\bfseries, text=mLightBrown, above=0.15cm of c3] {W (\texttt{hi})};
    \node[font=\scriptsize\bfseries, text=mDarkTeal, above=0.15cm of c0] {R (\texttt{lo})};
\end{tikzpicture}
\end{center}

\begin{itemize}
    \item O produtor escreve na posição \texttt{hi}; o consumidor lê da posição \texttt{lo}.
    \item Os índices andam com \texttt{(i + 1) \% N}: depois da última posição, voltam para a 0 (\textbf{buffer circular}).
    \item Buffer \textbf{cheio}: o produtor precisa esperar. Buffer \textbf{vazio}: o consumidor precisa esperar.
    \item Ninguém pode mexer no buffer enquanto o outro mexe.
\end{itemize}
\end{frame}

\begin{frame}[fragile]{Produtor e consumidor: \texttt{prod\_cons.c}}
\begin{columns}[T]
\column{0.5\textwidth}
\begin{lstlisting}[style=codigo, basicstyle=\ttfamily\scriptsize\color{termfg}]
/* PRODUTOR */
wait_sem(&empty);   /* 1. reserva VAGA  */
wait_sem(&mutex);   /* 2. pega a chave  */
insert_item(item);  /* 3. mexe no buffer */
sem_post(&mutex);   /* 4. devolve chave */
sem_post(&full);    /* 5. avisa: ITEM   */
\end{lstlisting}

\column{0.5\textwidth}
\begin{lstlisting}[style=codigo, basicstyle=\ttfamily\scriptsize\color{termfg}]
/* CONSUMIDOR */
wait_sem(&full);    /* 1. reserva ITEM  */
wait_sem(&mutex);   /* 2. pega a chave  */
item = remove_item(); /* 3. mexe       */
sem_post(&mutex);   /* 4. devolve chave */
sem_post(&empty);   /* 5. avisa: VAGA   */
\end{lstlisting}
\end{columns}

\vspace{0.4em}
\begin{exampleblock}{Regra para lembrar}
Cada thread \textbf{espera o que consome} e \textbf{avisa o que cria}. O produtor consome vagas e cria itens; o consumidor, o contrário.
\end{exampleblock}

\begin{lstlisting}[style=terminal, basicstyle=\ttfamily\scriptsize\color{termfg}]
$ gcc -Wall -Wextra examples/dia07/prod_cons.c -o prod_cons -pthread && ./prod_cons
Produzido: 29 | ocupacao: 1/20   ...   Fim: 40 itens produzidos e consumidos; buffer vazio.
\end{lstlisting}
\end{frame}

\begin{frame}[fragile]{E se a ordem for invertida?}
\begin{columns}[c]
\column{0.45\textwidth}
\begin{lstlisting}[style=codigo, basicstyle=\ttfamily\scriptsize\color{termfg}]
/* CONSUMIDOR ERRADO */
wait_sem(&mutex);   /* pega a chave...    */
wait_sem(&full);    /* ...e dorme com ela */
\end{lstlisting}

\small Com o buffer \textbf{vazio}:
\begin{enumerate}
    \item O consumidor pega o \texttt{mutex}.
    \item Dorme em \texttt{full}, \textbf{segurando a chave}.
    \item O produtor dorme em \texttt{mutex}.
\end{enumerate}

\column{0.53\textwidth}
\begin{center}
\begin{tikzpicture}[node distance=2.2cm]
    \node[proc] (c) {Consumidor\\\scriptsize tem \texttt{mutex}};
    \node[proc, right=of c] (p) {Produtor\\\scriptsize quer \texttt{mutex}};
    \draw[setad, bend left=25] (c) to node[rotulo, above] {espera um item} (p);
    \draw[setad, bend left=25] (p) to node[rotulo, below] {espera a chave} (c);
\end{tikzpicture}
\end{center}

\begin{alertblock}{Deadlock}
Cada um espera algo que só o outro pode dar. O programa trava para sempre.
\end{alertblock}
\end{columns}
\end{frame}

% ============================================================
\section{Ver funcionando: simulação}
% ============================================================

\begin{frame}[fragile]{Simulação ``Produção de dinheiro''}
\begin{columns}[T]
\column{0.5\textwidth}
\small
\begin{tabular}{@{}ll@{}}
\toprule
\textbf{Na tela} & \textbf{No \texttt{prod\_cons.c}} \\
\midrule
10 posições & \texttt{N} \\
W: próxima escrita & \texttt{hi} \\
R: próxima leitura & \texttt{lo} \\
itens disponíveis & \texttt{full} \\
vagas livres & \texttt{empty} \\
``Aguardando\ldots'' & dormindo num semáforo \\
\bottomrule
\end{tabular}

\column{0.48\textwidth}
\small
\begin{itemize}
    \item \textbf{Banco}: o produtor (cria notas).
    \item \textbf{Carro-forte}: o consumidor (retira notas).
    \item Por baixo: threads e três semáforos, como no C.
\end{itemize}

\begin{exampleblock}{Só para ver}
Cinco minutos para \textbf{enxergar} o que estudamos. Sem atividade.
\end{exampleblock}
\end{columns}

\begin{lstlisting}[style=terminal, basicstyle=\ttfamily\scriptsize\color{termfg}]
# compila e abre a janela (atalho: make run-simulation)
$ mvn -f pc_trabalho04_202011393/pom.xml javafx:run
\end{lstlisting}
\end{frame}

\begin{frame}{O que observar}
\begin{columns}[T]
\column{0.48\textwidth}
\textbf{Produtor 10 · Consumidor 0,2}

\vspace{0.4em}
O buffer enche: \textbf{``Produtor: Aguardando vaga''}.

\vspace{0.4em}
\begin{alertblock}{No C}
Dormindo em \texttt{sem\_wait(\&empty)}, com \texttt{empty = 0}, \textbf{sem} segurar o \texttt{mutex}.
\end{alertblock}

\column{0.48\textwidth}
\pause
\textbf{Produtor 0,2 · Consumidor 10}

\vspace{0.4em}
O buffer esvazia: \textbf{``Consumidor: Aguardando item''}.

\vspace{0.4em}
\begin{alertblock}{No C}
Dormindo em \texttt{sem\_wait(\&full)}, com \texttt{full = 0}, até o próximo \texttt{sem\_post(\&full)}.
\end{alertblock}
\end{columns}

\pause
\vspace{0.8em}
\begin{exampleblock}{Em uma frase}
O \texttt{mutex} garante que ninguém atrapalha. \texttt{empty} e \texttt{full} garantem que ninguém trabalha sem ter o que precisa.
\end{exampleblock}
\end{frame}

% ============================================================
\section{Por que monitores?}
% ============================================================

\begin{frame}{Por que monitores?}
Semáforos funcionam, mas deixam tudo nas mãos do programador:

\begin{itemize}
    \item É fácil \textbf{esquecer} um \texttt{sem\_post}: alguém dorme para sempre.
    \item É fácil \textbf{trocar a ordem}: deadlock, como acabamos de ver.
    \item É fácil \textbf{proteger a coisa errada}: a corrida continua.
\end{itemize}

\vspace{0.6em}
\begin{alertblock}{Monitor}
Um objeto que junta os \textbf{dados compartilhados}, as \textbf{operações} sobre eles e a \textbf{exclusão mútua}. Quem usa o objeto não precisa lembrar de pegar e devolver a chave.
\end{alertblock}

\small Analogia: um banheiro com uma \textbf{chave única}. Quem tem a chave entra; os outros esperam na porta.
\end{frame}

\begin{frame}{As três ferramentas de um monitor Java}
\begin{center}
\begin{tabular}{@{}>{\ttfamily}lR{9.6cm}@{}}
\toprule
\textrm{\textbf{Recurso}} & \textbf{O que faz} \\
\midrule
synchronized & Para executar o método, a thread precisa da \textbf{chave} (lock) do objeto. Só uma thread por vez dentro dos métodos \texttt{synchronized} do mesmo objeto. \\
wait() & \textbf{Solta a chave} e dorme até ser avisada. Só pode ser chamado com a chave na mão. \\
notify() & Acorda \textbf{uma} thread que esperava nesse objeto. Ela ainda precisa \textbf{pegar a chave de novo} para continuar. \\
notifyAll() & Acorda \textbf{todas}; cada uma confere a sua condição. \\
\bottomrule
\end{tabular}
\end{center}

\vspace{0.4em}
\small \texttt{Thread.sleep()} \textbf{não} solta a chave. Só \texttt{wait()} solta.
\end{frame}

\begin{frame}[standout]
Parte 2: prática

\vspace{0.6em}
{\normalsize Em duplas, no computador, com a folha de atividade}
\end{frame}

% ============================================================
\section{Prática 1: seis lacunas em C}
% ============================================================

\begin{frame}[fragile]{Seis lacunas: \texttt{prod\_cons\_atividade.c}}
\begin{lstlisting}[style=codigo, basicstyle=\ttfamily\scriptsize\color{termfg}]
sem_init(&mutex, 0, 1);
sem_init(&empty, 0, @/* 1: vagas no inicio */@);
sem_init(&full,  0, @/* 2: itens no inicio */@);
\end{lstlisting}

\begin{columns}[T]
\column{0.48\textwidth}
\begin{lstlisting}[style=codigo, basicstyle=\ttfamily\scriptsize\color{termfg}]
/* produtor */
wait_sem(@/* 3 */@);
wait_sem(&mutex);
insert_item(item);
sem_post(&mutex);
sem_post(@/* 4 */@);
\end{lstlisting}
\column{0.48\textwidth}
\begin{lstlisting}[style=codigo, basicstyle=\ttfamily\scriptsize\color{termfg}]
/* consumidor */
wait_sem(@/* 5 */@);
wait_sem(&mutex);
item = remove_item();
sem_post(&mutex);
sem_post(@/* 6 */@);
\end{lstlisting}
\end{columns}

\footnotesize
\textbf{1 e 2:} quantas vagas e quantos itens no início?\\
\textbf{3 e 5:} o que cada um precisa antes de mexer?\\
\textbf{4 e 6:} o que cada um avisa depois?\\
O \texttt{mutex} já está pronto. Não mude a ordem das linhas.
\end{frame}

\begin{frame}[fragile]{Compilar, testar e diagnosticar}
\begin{lstlisting}[style=terminal, basicstyle=\ttfamily\scriptsize\color{termfg}]
$ gcc -std=c11 -Wall -Wextra -Wpedantic -Werror \
      examples/dia07/prod_cons_atividade.c -o prod-cons-aluno -pthread
$ ./prod-cons-aluno
...
Fim: 40 itens produzidos e consumidos; buffer vazio.
\end{lstlisting}

\begin{center}
\small
\begin{tabular}{@{}R{5.6cm}R{7cm}@{}}
\toprule
\textbf{Sintoma} & \textbf{O que significa} \\
\midrule
Trava sem terminar (\texttt{Ctrl + C}) & Alguém espera um semáforo que ninguém libera \\
\texttt{Assertion `count < N' failed} & O produtor inseriu sem haver vaga \\
\texttt{Assertion `count > 0' failed} & O consumidor retirou sem haver item \\
\bottomrule
\end{tabular}
\end{center}

\small Antes de preencher, o esqueleto \textbf{não compila}. É proposital.
\end{frame}

% ============================================================
\section{Prática 2: Mailbox sem sincronização}
% ============================================================

\begin{frame}[fragile]{Rodando a Mailbox 3b}
\small Uma \texttt{Mailbox3b} compartilhada por \textbf{um consumidor} e \textbf{dois produtores}: Dave diz ``Hello, world.'' e Bill diz ``Hot dog!''.

\textbf{Previsão:} as mensagens vão chegar certas?

\begin{lstlisting}[style=terminal, basicstyle=\ttfamily\scriptsize\color{termfg}]
$ javac -d build/java-3b examples/dia07/mailbox-3b/*.java
$ timeout --foreground 5 java -cp build/java-3b ThreadSync3b | grep "My name"
\end{lstlisting}

\begin{itemize}\small
    \item O programa nunca termina sozinho: o \texttt{timeout} encerra em 5 segundos.
    \item Rode também \textbf{sem} o \texttt{| grep}: o que são as outras linhas?
    \item Abra \texttt{Mailbox3b.java} e leia \texttt{storeMessage()}.
\end{itemize}
\end{frame}

\begin{frame}[fragile]{Mailbox 3b: o que aconteceu}
\begin{lstlisting}[style=java, basicstyle=\ttfamily\scriptsize\color{termfg}]
public void storeMessage(String msg) {         // dois produtores chamam ao mesmo tempo
    for (int i = 0; i < msg.length(); i++) {
        message[i] = msg.charAt(i);            // escreve letra por letra no MESMO array
        Thread.sleep((int) (MAXPROCESSTIME * Math.random()));
    }
    youHaveMail = true;
}
\end{lstlisting}

\textbf{Saída real:}
\begin{lstlisting}[style=terminal, basicstyle=\ttfamily\scriptsize\color{termfg}]
My name is, Dave. I say, Helloog!
My name is, Bill. I say, Hotlo, world.
My name is, Dave. I say, Hot dog!
\end{lstlisting}

\small As letras dos dois produtores se misturam no mesmo array: é a condição de corrida, agora em Java.
\end{frame}

\begin{frame}{Prática 2: responda (questões 3 a 5)}
\begin{enumerate}
    \setcounter{enumi}{2}
    \item Qual é o \textbf{recurso compartilhado} neste exemplo?
    \item Onde está a \textbf{região crítica} no \texttt{Mailbox3b}?
    \item Por que aumentar \texttt{MAXPROCESSTIME} tende a tornar o problema \textbf{mais visível}?
\end{enumerate}

\vspace{1em}
\begin{exampleblock}{Se sobrar tempo}
Mude \texttt{MAXPROCESSTIME} de 7 para 50 em \texttt{Mailbox3b.java}, compile e rode de novo. Compare.
\end{exampleblock}
\end{frame}

% ============================================================
\section{Prática 3: Mailbox com monitor}
% ============================================================

\begin{frame}[fragile]{Mailbox 3s: com monitor}
\begin{lstlisting}[style=java, basicstyle=\ttfamily\scriptsize\color{termfg}]
public synchronized void storeMessage(String msg) {   // precisa da chave do objeto
    ... copia a mensagem letra por letra ...
    youHaveMail = true;
    notify();                                           // acorda o consumidor
}

public synchronized String retrieveMessage() {
    while (youHaveMail == false)
        wait();                                         // solta a chave e dorme
    String msg = new String(message, 0, MAXMESSAGE);
    youHaveMail = false;
    return msg;
}
\end{lstlisting}

\small Agora um produtor só começa a escrever quando o outro terminou: as mensagens saem \textbf{íntegras}. O consumidor dorme em vez de perguntar repetidamente.
\end{frame}

\begin{frame}[fragile]{Rodando a Mailbox 3s}
\textbf{Previsão:} as mensagens ainda vão se misturar? E o ``How sad, no mail''?

\begin{lstlisting}[style=terminal, basicstyle=\ttfamily\scriptsize\color{termfg}]
$ javac -d build/java-3s examples/dia07/mailbox-3s/*.java
$ timeout 5 java -cp build/java-3s ThreadSync3s
Looking for my mail ..
My name is, Dave. I say, Hello, world.
Looking for my mail ..
My name is, Bill. I say, Hot dog!
\end{lstlisting}

\begin{columns}[T]
\column{0.48\textwidth}
\small
\textbf{Versão 3b}: acesso concorrente sem coordenação; o consumidor pergunta e volta de mãos vazias.
\column{0.48\textwidth}
\small
\textbf{Versão 3s}: \texttt{synchronized} protege o estado; \texttt{wait()}/\texttt{notify()} fazem o consumidor dormir até ter mensagem.
\end{columns}

\vspace{0.6em}
\small \textbf{Questão 6:} o que deixou de aparecer em relação à 3b? Qual parte do código causa cada mudança?
\end{frame}

\begin{frame}{Por que \texttt{while} e não \texttt{if} antes de \texttt{wait()}?}
\begin{itemize}
    \item Quem acorda \textbf{não recebe a chave na hora}: precisa disputá-la de novo.
    \item Enquanto isso, outra thread pode entrar e \textbf{mudar a condição} (por exemplo, pegar a mensagem antes).
    \item O Java também permite \textbf{despertares espúrios}: acordar sem ninguém ter chamado \texttt{notify}.
\end{itemize}

\vspace{0.6em}
\begin{exampleblock}{Regra}
Ao acordar, \textbf{confira de novo}. Com \texttt{while}, a thread só segue quando a condição for verdadeira de fato.
\end{exampleblock}
\end{frame}

\begin{frame}[fragile]{O limite da mailbox 3s}
\begin{alertblock}{Ela ainda pode perder mensagens}
\texttt{storeMessage} não espera a caixa esvaziar. Se dois produtores escrevem antes de o consumidor ler, a primeira mensagem é \textbf{sobrescrita}. A 3s garante \textbf{integridade}, não \textbf{entrega}.
\end{alertblock}

\textbf{Como seria a versão completa} (uma posição, sem perda):
\begin{lstlisting}[style=terminal, basicstyle=\ttfamily\scriptsize\color{termfg}]
guardar (synchronized):                retirar (synchronized):
    enquanto cheia: wait()                 enquanto vazia: wait()
    copia a mensagem                       copia a mensagem
    cheia = verdadeiro                     cheia = falso
    notifyAll()                            notifyAll()
\end{lstlisting}

\small É o produtor-consumidor com N = 1, escrito com monitor em vez de semáforos.
\end{frame}

% ============================================================
\section{Atividade final}
% ============================================================

\begin{frame}{Atividade final (individual)}
\begin{enumerate}
    \setcounter{enumi}{6}
    \item O que um monitor acrescenta em relação à ideia de apenas compartilhar um objeto?
    \item Qual é o papel de \texttt{synchronized} no exemplo \texttt{Mailbox3s}?
    \item Por que \texttt{retrieveMessage()} usa \texttt{while} antes de \texttt{wait()} em vez de simplesmente continuar?
    \item Qual a diferença conceitual entre espera ocupada e \texttt{wait()}?
\end{enumerate}
\end{frame}

% ============================================================
\section{Fechamento}
% ============================================================

\begin{frame}{Erros comuns}
\begin{center}
\small
\begin{tabular}{@{}R{4.6cm}R{8.4cm}@{}}
\toprule
\textbf{Se você pensou\ldots} & \textbf{Reveja\ldots} \\
\midrule
``\texttt{contador++} é uma operação só.'' & São três passos: ler, somar, escrever. \\
``\texttt{volatile} resolve a corrida.'' & Não: é preciso exclusão mútua. \\
``O \texttt{mutex} conta as vagas.'' & O \texttt{mutex} é a chave; quem conta é \texttt{empty} e \texttt{full}. \\
``Tanto faz a ordem dos \texttt{sem\_wait}.'' & Dormir segurando a chave causa deadlock. \\
``Semáforo é igual a Peterson.'' & No semáforo quem espera dorme; em Peterson, gira. \\
``\texttt{notify} entrega a chave na hora.'' & Quem acorda disputa a chave e confere de novo. \\
``\texttt{sleep} libera o monitor.'' & Só \texttt{wait} libera. \\
``A mailbox 3s entrega tudo.'' & Ela pode sobrescrever mensagem não lida. \\
\bottomrule
\end{tabular}
\end{center}
\end{frame}

\begin{frame}{Síntese}
\begin{center}
\begin{tabular}{@{}lR{9.6cm}@{}}
\toprule
\textbf{Condição de corrida} & Resultado depende da ordem de intercalação das threads. \\
\textbf{Região crítica} & Trecho que acessa o recurso compartilhado. \\
\textbf{Peterson} & Exclusão mútua para duas threads, com espera ocupada. \\
\textbf{Semáforo} & Contador com \texttt{wait}/\texttt{post}; quem espera dorme. \\
\textbf{Produtor-consumidor} & \texttt{mutex} para exclusão; \texttt{empty} e \texttt{full} para disponibilidade. \\
\textbf{Monitor} & Dados, operações e exclusão mútua juntos; \texttt{wait}/\texttt{notify} para condições. \\
\bottomrule
\end{tabular}
\end{center}

\vspace{0.6em}
\begin{exampleblock}{As duas aulas em uma frase}
Processos isolam a memória e precisam de pipe para conversar. Threads compartilham e precisam de sincronização para não se atrapalhar.
\end{exampleblock}
\end{frame}

\begin{frame}{Glossário}
\footnotesize
\begin{columns}[T]
\column{0.5\textwidth}
\textbf{Thread}: fluxo de execução dentro de um processo; divide a memória com as outras.\\[0.35em]
\textbf{Condição de corrida}: resultado depende da ordem das threads.\\[0.35em]
\textbf{Data race}: acesso concorrente sem sincronização; em C, comportamento indefinido.\\[0.35em]
\textbf{Região crítica}: trecho que mexe no recurso compartilhado.\\[0.35em]
\textbf{Exclusão mútua}: uma thread por vez na região crítica.\\[0.35em]
\textbf{Espera ocupada}: esperar girando e gastando CPU.

\column{0.5\textwidth}
\textbf{Semáforo}: contador não negativo com \texttt{wait} e \texttt{post}.\\[0.35em]
\textbf{Mutex}: semáforo (ou lock) que vale 0 ou 1; a ``chave''.\\[0.35em]
\textbf{Deadlock}: cada um espera algo que só o outro pode dar.\\[0.35em]
\textbf{Monitor}: objeto com dados, operações e exclusão mútua juntos.\\[0.35em]
\textbf{\texttt{wait} / \texttt{notify}}: dormir soltando a chave / acordar quem dorme.\\[0.35em]
\textbf{Buffer circular}: índices voltam a 0 com \texttt{\% N}.
\end{columns}
\end{frame}

\begin{frame}[fragile]{Cola: compilar e executar os exemplos}
\small Na raiz do repositório:
\begin{lstlisting}[style=terminal, basicstyle=\ttfamily\tiny\color{termfg}]
# C (sempre com -pthread; && so executa se a compilacao deu certo)
$ gcc -Wall -Wextra examples/dia07/corrida.c -o corrida -pthread && ./corrida
$ gcc -Wall -Wextra examples/dia07/corrida_atomic.c -o corrida-atomic -pthread && ./corrida-atomic
$ gcc -Wall -Wextra examples/dia07/corrida_mutex.c -o corrida-mutex -pthread && ./corrida-mutex
$ gcc -Wall -Wextra examples/dia07/peterson-code.c -o peterson -pthread && ./peterson
$ gcc -Wall -Wextra examples/dia07/prod_cons.c -o prod_cons -pthread && ./prod_cons

# Java (as mailboxes nunca terminam: timeout encerra em 5 s)
$ javac -d build/java-3b examples/dia07/mailbox-3b/*.java && timeout 5 java -cp build/java-3b ThreadSync3b
$ javac -d build/java-3s examples/dia07/mailbox-3s/*.java && timeout 5 java -cp build/java-3s ThreadSync3s
$ javac -d build/caixa examples/dia07/caixa/*.java && java -cp build/caixa TesteCaixa

# Simulacao Producao de dinheiro
$ mvn -f pc_trabalho04_202011393/pom.xml javafx:run

# Pratica 1 (avisos viram erros)
$ gcc -std=c11 -Wall -Wextra -Wpedantic -Werror \
      examples/dia07/prod_cons_atividade.c -o prod-cons-aluno -pthread && ./prod-cons-aluno
\end{lstlisting}
\end{frame}

\begin{frame}{Para estudar depois}
\textbf{Rode os exemplos}: os comandos estão no slide anterior.

\vspace{0.5em}
\textbf{Leia}
\begin{itemize}
    \item A \textbf{apostila} da aula: os mesmos temas, com mais calma, e exercícios com respostas.
    \item \texttt{man 3 sem\_wait}, \texttt{man 3 pthread\_create}, \texttt{man 7 pthreads}.
    \item Tanenbaum e Bos, \emph{Sistemas Operacionais Modernos}, 4. ed., seção 2.3 (comunicação entre processos).
    \item Silberschatz, Galvin e Gagne, \emph{Fundamentos de Sistemas Operacionais}, 9. ed., cap. 5 (sincronização de processos).
\end{itemize}
\end{frame}

\begin{frame}{Obrigado!}
\centering
{\Large Dúvidas?}
\vspace{2em}

Weslei Ferreira Santos\\
MATA58 · Sistemas Operacionais\\
Instituto de Computação / UFBA

\vspace{1.5em}
{\scriptsize Material-base: Prof. Maycon Leone M. Peixoto}
\end{frame}

\end{document}
